\chapter{Ce que fait une société de trading}\label{ch:m1:what-a-trading-firm-does}

À 10~h~31~min~04~s, cent actions d'une grande entreprise changent de mains à 50.01~\$.
L'acheteuse est une dentiste de l'Ohio qui a appuyé sur un bouton de son téléphone~; le vendeur est
une société installée dans un centre de données, qui ne l'a jamais connue et détenait les actions depuis
onze secondes. Entre les deux parties, un courtier, une bourse et une chambre de compensation ont chacun
traité l'ordre, et chacun a été rémunéré. La transaction n'a duré qu'une microseconde sur le
tape. Ce chapitre porte sur ce que sont ces sociétés, sur la façon dont chacune tire ses revenus
de cette transaction, et sur la raison pour laquelle aucune d'elles n'exerce le même métier.

\section{Cinq modèles d'activité}

Toutes les sociétés de ce livre vivent du trading, mais l'expression «~société de trading~» recouvre au
moins cinq activités différentes. Elles se distinguent par trois éléments~: à qui appartient le capital
exposé au risque, qui les rémunère et combien de temps elles conservent une position.

\begin{definition}[Teneur de marché]\label{def:m1:what-a-trading-firm-does:market-maker}
Un \emph{teneur de marché}\index{teneur de marché} (market maker) est une société qui propose en permanence
d'acheter et de vendre un même instrument, à un prix d'achat (son
\emph{cours acheteur}) inférieur à son prix de vente (son \emph{cours vendeur}), et qui conclut avec quiconque accepte.
Elle risque son propre capital et vise à clore chaque journée avec une position faible,
voire nulle.
\end{definition}

\begin{definition}[Fourchette et prix milieu]\label{def:m1:what-a-trading-firm-does:spread}
Pour un meilleur cours acheteur $b$ et un meilleur cours vendeur $a > b$, la
\emph{fourchette}\index{fourchette} est $s = a - b$ et le
\emph{prix milieu}\index{prix milieu} est $m = \tfrac12(a+b)$. La
\emph{demi-fourchette} $s/2$ est ce que paie un ordre au marché par rapport au prix milieu.
\end{definition}

\begin{example}[L'ordre de la dentiste]\label{ex:m1:what-a-trading-firm-does:dentist}
Le marché affiche $49.99$ à l'achat et $50.01$ à la vente~: $s = \$0.02$, $m = \$50.00$, et la
\omterm{def:m1:what-a-trading-firm-does:spread}{fourchette} vaut $0.02/50.00 = \qty{4}{\bp}$ du prix. L'ordre au marché de la dentiste
achète au cours vendeur et paie la \omterm{def:m1:what-a-trading-firm-does:spread}{demi-fourchette}, soit 0.01~\$ par action, c'est-à-dire
1.00~\$ sur ses cent actions. Ce dollar est le revenu brut du \omterm{def:m1:what-a-trading-firm-does:market-maker}{teneur de marché}
sur la transaction, avant que le cours n'ait eu le temps de bouger.
\end{example}

\begin{definition}[Société de trading pour compte propre]\label{def:m1:what-a-trading-firm-does:prop}
Une \emph{société de trading pour compte propre}\index{société de trading pour compte propre} (proprietary trading firm) négocie
uniquement le capital de ses propriétaires, n'a pas de clients et ne perçoit aucune commission~:
son revenu tout entier se résume à son profit de trading. La plupart des \omterm{def:m1:what-a-trading-firm-does:market-maker}{teneurs de marché}
électroniques sont des sociétés pour compte propre~; beaucoup de ces sociétés mettent aussi en œuvre
des stratégies qui prennent de la liquidité au lieu d'en fournir.
\end{definition}

\begin{definition}[Hedge fund et plateforme multigérants]\label{def:m1:what-a-trading-firm-does:hedge-fund}
Un \emph{hedge fund}\index{hedge fund} est un véhicule d'investissement offert en privé, qui gère
le capital d'investisseurs extérieurs, avec peu de contraintes sur les instruments,
l'effet de levier ou la vente à découvert, et qui perçoit des frais sur les actifs, plus
une part des profits. Une \emph{plateforme multigérants}\index{plateforme multigérants}
est un hedge fund organisé en de nombreuses équipes indépendantes («~pods~»), chacune dotée
d'un capital et de limites de risque strictes accordées par un allocateur central qui
compense les positions et gère l'ensemble.
\end{definition}

\begin{definition}[Société de gestion et encours sous gestion]\label{def:m1:what-a-trading-firm-does:asset-manager}
Une \emph{société de gestion}\index{société de gestion} investit l'argent de ses clients dans le cadre
d'un mandat écrit, généralement pour répliquer ou pour surperformer un indice de référence, et perçoit
un pourcentage des \emph{encours sous gestion}\index{encours sous gestion}
(AUM), c'est-à-dire la valeur de marché des portefeuilles qu'elle gère.
\end{definition}

Le cinquième modèle est le desk de trading d'une banque~: un dealer qui affiche des prix à
ses clients en échange de leur activité, étudié au \cref{ch:m1:the-sell-side}.
L'\Cref{ex:m1:what-a-trading-firm-does:table} présente les cinq modèles côte à côte.

\begin{example}[Cinq sociétés, un tableau]\label{ex:m1:what-a-trading-firm-does:table}
\begin{center}\footnotesize
\begin{tabular}{@{}lllll@{}}
\hline
 & Capital de & Rémunération & Détention & Coût dominant \\
\hline
\omterm{def:m1:what-a-trading-firm-does:market-maker}{Teneur de marché} & propriétaires & la \omterm{def:m1:what-a-trading-firm-does:spread}{fourchette} & secondes à heures & technologie \\
Société pour compte propre & propriétaires & profit de trading & secondes à semaines & technologie, personnel \\
\omterm{def:m1:what-a-trading-firm-does:hedge-fund}{Hedge fund} & investisseurs & frais (actifs, profits) & jours à années & personnel \\
\omterm{def:m1:what-a-trading-firm-does:asset-manager}{Société de gestion} & clients & frais sur actifs & mois à années & distribution \\
Desk de banque & actionnaires & \omterm{def:m1:what-a-trading-firm-does:spread}{fourchette}, clientèle & minutes à mois & capital \\
\hline
\end{tabular}
\end{center}
La suite de la série revient sans cesse à ce tableau~: une technique qui est centrale pour une ligne
(la vitesse, pour la première) est sans intérêt pour une autre (la
quatrième).
\end{example}

\begin{remark}[D'où vient le «~two and twenty~»]\label{rem:m1:what-a-trading-firm-does:jones}
La structure de frais des \omterm{def:m1:what-a-trading-firm-does:hedge-fund}{hedge funds} est aussi ancienne que le \omterm{def:m1:what-a-trading-firm-does:hedge-fund}{hedge fund} lui-même. Alfred Winslow
Jones, qui lança le premier en 1949 avec 100\,000~\$, prélevait 20~\% des
profits et, au début, rien d'autre~; des frais annuels sur les actifs vinrent plus tard, et
«~2 and 20~» --- 2~\% des actifs, 20~\% des profits --- devint la référence à partir de laquelle
chaque fonds négocie aujourd'hui à la baisse ou, pour les plus recherchés,
à la hausse.
\end{remark}

\begin{dated}{2026-09}{Taille de l'industrie}\label{dat:m1:what-a-trading-firm-does:aum}
Les actifs gérés professionnellement dans le monde s'élevaient à environ 147~billions de dollars à la fin
de 2025, contre 128~billions un an plus tôt, selon le décompte du secteur le plus souvent cité.
L'essentiel de la croissance de cette année-là provient de la hausse des marchés, et non de nouveaux
apports d'argent.
\end{dated}

\section{D'où vient l'argent}

Une fois le vocabulaire mis de côté, une activité de trading n'a que cinq sources de
revenus.

\begin{method}[Classer un revenu de trading]\label{met:m1:what-a-trading-firm-does:sources}
Pour chaque dollar de revenu de trading, déterminer à laquelle de ces sources il appartient~:
\begin{enumerate}
\item \textbf{\omterm{def:m1:what-a-trading-firm-does:spread}{Fourchette}~:} payée par quelqu'un qui voulait échanger \emph{maintenant}.
\item \textbf{Frais ou commission~:} payés pour un service (gestion, exécution, compensation, prêt,
  diffusion de données) quel que soit le comportement du marché.
\item \textbf{Prime de risque~:} payée, en moyenne, pour détenir un risque que d'autres veulent céder
  --- risque sur actions, manque de liquidité, côté vendeur d'une assurance.
\item \textbf{Carry et financement~:} la différence entre le rendement d'une position
  et son coût de financement.
\item \textbf{\omterm{def:m1:what-a-trading-firm-does:alpha}{Alpha}~:} profit tiré d'un pari juste sur un prix alors que d'autres
  se trompaient.
\end{enumerate}
Un \omterm{def:m1:what-a-trading-firm-does:market-maker}{teneur de marché} vit de la source~1~; une \omterm{def:m1:what-a-trading-firm-does:asset-manager}{société de gestion} vit de la source~2~; un investisseur indiciel perçoit la source~3~;
le desk de financement d'une banque vit de la source~4~; les \omterm{def:m1:what-a-trading-firm-does:hedge-fund}{hedge funds} vendent la source~5 et livrent souvent
un mélange des sources~3 et~5.
\end{method}

\begin{definition}[Alpha]\label{def:m1:what-a-trading-firm-does:alpha}
L'\emph{alpha}\index{alpha} d'une stratégie est la part de son rendement espéré
qui ne rémunère pas la prise de risques accessibles à bon marché
ailleurs~: dans une régression de ses rendements sur ceux des facteurs de
risque négociables, il s'agit de l'ordonnée à l'origine.
\end{definition}

La source~1 ressemble à de l'argent facile~: acheter au cours acheteur, vendre au cours vendeur, recommencer. Ce
n'est pas le cas, et la raison en est l'idée la plus importante du chapitre.

\begin{definition}[Sélection adverse]\label{def:m1:what-a-trading-firm-does:adverse-selection}
Un \omterm{def:m1:what-a-trading-firm-does:market-maker}{teneur de marché} subit une \emph{sélection adverse}\index{sélection adverse}
lorsque les contreparties qui choisissent de traiter avec ses cotations sont, plus
souvent que par hasard, celles qui savent que le prix est sur le point de se déplacer contre lui. Son coût
se mesure comme le mouvement moyen du \omterm{def:m1:what-a-trading-firm-does:spread}{prix milieu} contre la nouvelle position
du \omterm{def:m1:what-a-trading-firm-does:market-maker}{teneur de marché} après une transaction.
\end{definition}

\begin{proposition}[Captation nette]\label{prop:m1:what-a-trading-firm-does:capture}
Un \omterm{def:m1:what-a-trading-firm-does:market-maker}{teneur de marché} cote $m \pm s/2$. Une fraction $p$ des ordres qu'il exécute sont
informés~: après chacun, le \omterm{def:m1:what-a-trading-firm-does:spread}{prix milieu} se déplace de $J$ contre le \omterm{def:m1:what-a-trading-firm-does:market-maker}{teneur de marché}.
Les autres ordres ne portent aucune information. Avec une rémunération $\rho$ et des frais
de compensation $\kappa$ par action, le profit espéré par action négociée, la
\emph{captation nette}\index{captation nette}, vaut
\[
c \;=\; \frac{s}{2} \;-\; pJ \;+\; \rho \;-\; \kappa .
\]
\end{proposition}
\begin{proof}
Un achat client d'une action à $m + s/2$ laisse le \omterm{def:m1:what-a-trading-firm-does:market-maker}{teneur de marché} vendeur
d'une action à ce prix. Valorisée au \omterm{def:m1:what-a-trading-firm-does:spread}{prix milieu}, la position a gagné $s/2$. Avec
probabilité $p$, le \omterm{def:m1:what-a-trading-firm-does:spread}{prix milieu} monte ensuite de $J$, et la position vendeuse perd $J$~; sinon,
le mouvement espéré est nul. Les ventes sont symétriques. La rémunération et les frais sont
payés par action, quoi qu'il arrive ensuite. Les positions héritées de transactions antérieures ont un
profit espéré nul, car le sens du prochain ordre est indépendant de
ces positions.
\end{proof}

\begin{example}[Un centime n'est pas un centime]\label{ex:m1:what-a-trading-firm-does:capture}
Prenons $s/2 = 1.00$ centime, $p = 15\%$, $J = 3$ centimes, $\rho = 0.20$ centime et
$\kappa = 0.02$ centime. Alors $c = 1.00 - 0.45 + 0.20 - 0.02 = 0.73$ centime~: près
de la moitié de la \omterm{def:m1:what-a-trading-firm-does:spread}{fourchette} gagnée est rendue à des contreparties mieux informées.
Si $p$ monte à $35\%$, $c = 0.13$ centime~; à $40\%$ l'activité perd de l'argent
sur chaque action tout en «~gagnant la \omterm{def:m1:what-a-trading-firm-does:spread}{fourchette}~» sur chacune.
\end{example}

\begin{proposition}[Volume du point mort]\label{prop:m1:what-a-trading-firm-does:breakeven}
Une société dont les coûts fixes s'élèvent à $C$ par jour et dont la \omterm{prop:m1:what-a-trading-firm-does:capture}{captation
nette} vaut $c > 0$ par action atteint le point mort à un volume quotidien $V^\star = C/c$,
et son profit quotidien au volume $V$ vaut $c\,(V - V^\star)$.
\end{proposition}
\begin{proof}
Le profit quotidien vaut $cV - C = c\,(V - C/c)$.
\end{proof}

La proposition est triviale, et elle résume toute l'économie de la \omterm{def:m1:what-a-trading-firm-does:market-maker}{teneur
de marché} électronique~: les coûts sont fixes (personnel, machines, données, connectivité) et
le revenu est un tout petit nombre multiplié par un très grand. Au-delà de $V^\star$
presque chaque action supplémentaire rapporte un profit~; en deçà, la société perd de l'argent. C'est pourquoi le secteur se
concentre, et pourquoi les résultats des \omterm{def:m1:what-a-trading-firm-does:market-maker}{teneurs de marché} cotés
suivent les volumes de marché.

\begin{omfigure}
\begin{tikzpicture}
\begin{axis}[width=11cm,height=5cm,
  xlabel={volume quotidien (millions d'actions)},ylabel={profit annuel (millions de \$)},
  tick label style={font=\small},label style={font=\small},
  xmin=0,xmax=60,ymin=-30,ymax=50,grid=major,grid style={gray!25}]
\addplot[omDef,very thick] table[col sep=comma,x=volume_m,y=profit_m]{figdata/markets-1/01-what-a-trading-firm-does/break_even.csv};
\addplot[omThm,dashed] coordinates {(20,-30) (20,50)};
\node[omThm,font=\small,anchor=west] at (axis cs:21,-10) {$V^\star = 20$ millions};
\end{axis}
\end{tikzpicture}

\omcaption{Profit annuel en fonction du volume quotidien pour une \omterm{prop:m1:what-a-trading-firm-does:capture}{captation nette}
de 0.45 centime par action et des coûts fixes de 90\,000~\$ par jour. En deçà de $V^\star$, la société perd
jusqu'à la totalité de sa base de coûts~; au-delà, le profit croît sans nouveau coût. Données~:
calculées par le script du chapitre.}
\end{omfigure}

\begin{dated}{2026-09}{Deux teneurs de marché cotés, d'après leurs propres chiffres}\label{dat:m1:what-a-trading-firm-does:listed}
La plupart des \omterm{def:m1:what-a-trading-firm-does:market-maker}{teneurs de marché} sont privés, mais deux grands d'entre eux publient leurs résultats.
Virtu Financial a déclaré un chiffre d'affaires total de 3\,632~millions de dollars et un résultat net
de négociation ajusté de 2\,145~millions de dollars pour 2025 (1\,598~millions de dollars pour 2024).
Flow Traders, spécialiste des produits cotés en bourse, a déclaré un résultat net
de négociation de 486 millions d'euros pour 2025, avec 635 employés en fin d'année~; pour 2024
il a déclaré 468 millions d'euros, 609 employés et 1\,545 milliards d'euros de
produits cotés en bourse négociés. L'exercice~\ref{exo:m1:what-a-trading-firm-does:6}
le transforme en une captation exprimée en points de base.
\end{dated}

\section{Principal et agent}

\begin{definition}[Principal et agent]\label{def:m1:what-a-trading-firm-does:principal-agent}
Une société négocie en \emph{principal}\index{principal} lorsqu'elle est elle-même
l'acheteuse ou la vendeuse~: la position, et son risque, entrent dans son propre
bilan. Elle négocie en \emph{agent}\index{agent} lorsqu'elle organise une
transaction entre son client et un tiers~: elle ne détient jamais la position et
touche une commission.
\end{definition}

\begin{example}[La même vente, deux fois]\label{ex:m1:what-a-trading-firm-does:twice}
Un client veut vendre 500\,000 actions cotées autour de 50~\$. En tant que
\emph{\omterm{def:m1:what-a-trading-firm-does:principal-agent}{agent}}, le courtier exécute l'ordre sur le marché au cours de la journée, pour une
commission de 1 centime par action~: il touche 5\,000~\$ quel que soit le prix
obtenu par le client, et c'est le client qui supporte le risque que le cours baisse pendant la journée.
En tant que \emph{\omterm{def:m1:what-a-trading-firm-does:principal-agent}{principal}}, le dealer achète tout le bloc sur-le-champ à 49.80~\$~:
le client est débarrassé de son ordre, et le dealer détient 24.9~millions de dollars d'actions qu'il doit
revendre. S'il les revend en moyenne à 49.92~\$, il gagne 60\,000~\$~; à 49.70~\$, il
perd 50\,000~\$. La décote de 20 centimes est le prix du transfert du risque.
\end{example}

\begin{remark}[Pourquoi la distinction est contrôlée]\label{rem:m1:what-a-trading-firm-does:conflict}
Un \omterm{def:m1:what-a-trading-firm-does:principal-agent}{agent} doit à son client le meilleur résultat qu'il peut obtenir~; un \omterm{def:m1:what-a-trading-firm-does:principal-agent}{principal} est
la contrepartie de son client et tire profit d'un prix plus défavorable pour lui.
Une société qui fait les deux --- la plupart des banques et des partenaires d'exécution des courtiers de détail ---
doit indiquer au client à quel titre elle a agi sur chaque transaction, et les régulateurs
des deux côtés de l'Atlantique examinent la manière dont elle a choisi. Les règles font
l'objet des \cref{ch:m1:retail-flow-and-wholesaling,ch:m1:european-equity-market-structure}.
\end{remark}

\section{Carte d'une transaction}

\begin{omfigure}
\begin{tikzpicture}[font=\small,
  box/.style={draw=omDef,rounded corners=2pt,minimum width=2.5cm,minimum height=0.9cm,align=center,fill=omDef!4},
  svc/.style={draw=omExo,rounded corners=2pt,minimum width=2.5cm,minimum height=0.9cm,align=center,fill=omExo!6},
  flow/.style={-{Stealth[length=2mm]},thick,omDef},
  pay/.style={-{Stealth[length=2mm]},dashed,omThm}]
\node[box] (inv) at (0,0) {Investisseur};
\node[box] (brk) at (4.3,0) {Courtier\\(agent)};
\node[box] (exc) at (8.6,0) {Bourse};
\node[box] (mm) at (12.9,0) {Teneur\\de marché\\(principal)};
\node[svc] (ccp) at (8.6,-2.6) {Chambre de compensation};
\node[svc] (dat) at (8.6,2.4) {Fournisseurs de données};
\draw[flow] (inv) -- node[above]{achat} (brk);
\draw[flow] (brk) -- node[above]{ordre} (exc);
\draw[flow] (mm) -- node[above]{offre} (exc);
\draw[flow] (exc) -- node[right]{transaction} (ccp);
\draw[flow] (exc) -- node[right]{cours} (dat);
\draw[pay] (inv.south) to[bend right=25] node[below,pos=0.5]{commission} (brk.south);
\draw[pay] (brk.south east) to[bend right=20] node[below left,pos=0.75]{frais du preneur} (exc.south west);
\draw[pay] (exc.south east) to[bend right=20] node[below right,pos=0.3]{rebate} (mm.south west);
\draw[pay] (mm.south) to[bend left=25] node[below right,pos=0.5,align=left]{frais de\\compensation} (ccp.east);
\end{tikzpicture}

\omcaption{Un ordre d'achat de détail. Les flèches pleines transportent les ordres et l'information~;
les flèches en pointillés transportent des sommes autres que le cours des actions. La \omterm{def:m1:what-a-trading-firm-does:spread}{demi-fourchette}
n'apparaît pas comme une flèche~: elle est incluse dans le cours que paie l'investisseur.}
\end{omfigure}

Suivons les cent actions du dentiste. Son courtier, agissant en \omterm{def:m1:what-a-trading-firm-does:principal-agent}{agent}, achemine
l'ordre vers une bourse et facture une commission (souvent nulle pour elle en apparence, et
récupérée ailleurs, comme on le verra au \cref{ch:m1:retail-flow-and-wholesaling}).
La bourse la confronte au cours vendeur affiché par le \omterm{def:m1:what-a-trading-firm-does:market-maker}{teneur de marché}~; selon le
barème \emph{maker--taker} courant, elle facture le côté qui a pris la liquidité
et reverse une partie de ce montant au côté qui l'a fournie
(exposé au \cref{ch:m1:us-equity-market-structure}). La chambre de compensation s'interpose entre
les deux côtés et garantit le règlement-livraison, moyennant des frais
(traité au \cref{ch:m1:clearing-and-settlement}). La bourse vend les données relatives à la
transaction aux fournisseurs de données et directement aux sociétés
(exposé au \cref{ch:m1:exchanges-brokers-venues}). Et le \omterm{def:m1:what-a-trading-firm-does:market-maker}{teneur de marché}, qui se trouve à découvert
de cent actions vendues à 50.01~\$, attend l'arrivée d'un vendeur à 49.99~\$ ---
et espère que le dentiste ne savait rien.

\section{Tutoriel~: qui a gagné quoi aujourd'hui~?}\label{tut:m1:what-a-trading-firm-does:whoearns}

\begin{tutorial}
\textbf{Objectif.} Simuler une journée de négociation sur un titre et répartir chaque
dollar échangé entre le \omterm{def:m1:what-a-trading-firm-does:market-maker}{teneur de marché}, la bourse, le courtier
et la chambre de compensation. \textbf{État final~:} le diagramme en barres qui suit ce tutoriel,
et une vérification de
la \cref{prop:m1:what-a-trading-firm-does:capture} à l'aide de la simulation.
\begin{tutsteps}
\item \textbf{Lire le modèle.} Vingt mille ordres clients de 100 actions
frappent un cours affiché de $m \pm 1$ cent~; 15~\% d'entre eux sont informés et font bouger le \omterm{def:m1:what-a-trading-firm-does:spread}{prix milieu} de
3 cents. Les frais suivent un barème maker--taker.
\omcode{code/markets-1/01-what-a-trading-firm-does/python/whoearns.py}{30}{41}{La boucle de négociation~: le teneur de marché prend l'autre côté de chaque ordre, et les informés font ensuite bouger le prix.}\label{lst:m1:what-a-trading-firm-does:loop}
\item \textbf{Décomposer le résultat du \omterm{def:m1:what-a-trading-firm-does:market-maker}{teneur de marché}.} La position est valorisée au
\omterm{def:m1:what-a-trading-firm-does:spread}{prix milieu} de clôture. Ce qui n'est pas gagné sur la \omterm{def:m1:what-a-trading-firm-does:spread}{fourchette} relève du gain ou de la perte
liée à la détention d'un stock pendant que les prix bougent~: surtout de la \omterm{def:m1:what-a-trading-firm-does:adverse-selection}{sélection adverse}.
\omcode{code/markets-1/01-what-a-trading-firm-does/python/whoearns.py}{42}{55}{Qui reçoit quoi. Le coût pour les clients est égal à la fourchette augmentée de la commission~; les professionnels se le partagent ensuite.}\label{lst:m1:what-a-trading-firm-does:split}
\item \textbf{Lancer la simulation.} Depuis une invite Python, dans le répertoire
\texttt{python/} du chapitre, importer \texttt{whoearns} et afficher
\texttt{simulate\_day(1)}. On doit obtenir une \omterm{def:m1:what-a-trading-firm-does:spread}{fourchette} gagnée de 20\,000~\$, une perte de position proche de 5\,600~\$
et un résultat net du \omterm{def:m1:what-a-trading-firm-does:market-maker}{teneur de marché} proche de 18\,000~\$.
\item \textbf{Comparer avec la théorie.} La proposition prévoit une perte de position
de $pJ = 0.45$ cent par action, soit 9\,000~\$ sur deux millions d'actions. Une journée est un
tirage~: en lançant mille graines, la moyenne converge vers elle, avec un écart-type
d'environ 17\,000~\$ par jour (l'histogramme ci-dessous).
\end{tutsteps}
\textbf{Pour aller plus loin.} Fixer \texttt{informed=0.35}~: la société
gagne toujours 20\,000~\$ de \omterm{def:m1:what-a-trading-firm-does:spread}{fourchette} par jour et arrive tout juste à l'équilibre. Puis faire
élargir après chaque transaction informée la cotation du \omterm{def:m1:what-a-trading-firm-does:market-maker}{teneur de marché}, et observer le volume ---
que le modèle maintient fixe --- devient ce qu'il faudrait modéliser.
\end{tutorial}

\begin{omfigure}
\begin{tikzpicture}
\begin{axis}[width=11.5cm,height=5.4cm,xbar,bar width=7pt,
  xlabel={milliers de dollars, une journée simulée},
  ytick=data,yticklabels={{Spread capté},{P\&L de position},{Rebates},{Net du teneur de marché},{Bourse},{Courtier},{Chambre de compensation}},
  yticklabel style={font=\small},tick label style={font=\small},label style={font=\small},
  y dir=reverse,xmin=-10,xmax=25,enlarge y limits=0.1,grid=major,grid style={gray!25},
  table/col sep=comma]
\addplot[fill=omDef!60,draw=omDef] table[x=usd_thousands,y=k]{figdata/markets-1/01-what-a-trading-firm-does/whoearns.csv};
\end{axis}
\end{tikzpicture}

\omcaption{Qui a gagné quoi lors d'une journée simulée portant sur deux millions d'actions. Le
résultat net du \omterm{def:m1:what-a-trading-firm-does:market-maker}{teneur de marché} correspond aux trois premières barres moins ses frais de compensation. Données~:
le simulateur du tutoriel, graine 1.}
\end{omfigure}

\begin{omfigure}
\begin{tikzpicture}
\begin{axis}[width=11cm,height=4.8cm,ybar,bar width=5.2pt,
  xlabel={net quotidien du teneur de marché (milliers de dollars)},ylabel={jours sur 1\,000},
  tick label style={font=\small},label style={font=\small},
  xmin=-62,xmax=82,ymin=0,xtick={-60,-40,-20,0,20,40,60,80},table/col sep=comma]
\addplot[fill=omDef!40,draw=omDef] table[x=centre,y=count]{figdata/markets-1/01-what-a-trading-firm-does/daily_net_hist.csv};
\end{axis}
\end{tikzpicture}

\omcaption{Mille journées simulées. La moyenne est la \omterm{prop:m1:what-a-trading-firm-does:capture}{captation nette} multipliée par le
volume~; la dispersion autour de celle-ci est le risque d'inventaire. Six journées sur sept sont
rentables. Données~: le simulateur du tutoriel, graines de 0 à 999.}
\end{omfigure}

\section{Construction~: le grand livre de la société}\label{bld:m1:what-a-trading-firm-does:ledger}

\begin{build}
\bfield{Rôle} La société de trading miniature assemblée tout au long de la série
a besoin d'un lieu unique où chaque dollar gagné ou dépensé est enregistré et
classé. Chaque composant ultérieur --- le gardien du P\&L
du \cref{ch:m1:pnl-and-positions}, les moteurs de frais, le calculateur de financement ---
y enregistre ses écritures.
\bfield{Interface} Un objet immuable \texttt{Entry(ts, account, category,
currency, amount, ref)} et un \texttt{Ledger} doté de \texttt{post(entry)},
\texttt{balance(account\_prefix, currency)} et
\texttt{by\_category(account\_prefix, currency)}. Les comptes sont des chemins pointés
(\texttt{desk.mm.equities})~; un préfixe sélectionne un sous-arbre. Les catégories sont
\texttt{spread}, \texttt{position}, \texttt{fee}, \texttt{rebate},
\texttt{commission}, \texttt{financing}, \texttt{other}.
\bfield{Règles} Les montants sont des \emph{entiers} signés, exprimés en dix-millièmes de
l'unité monétaire, de sorte qu'un rabais de 0.2 cent vaille exactement 20 unités~: jamais
d'arithmétique à virgule flottante pour les montants qui doivent se cumuler exactement. Les écritures sont enregistrées dans l'ordre chronologique~;
les catégories inconnues, les devises mal formées et les montants non entiers sont
rejetés. Les devises ne sont jamais additionnées entre elles.
\bfield{Tests d'acceptation} \texttt{code/firm/ledger/tests/test\_ledger.py}~:
dix mille rabais inférieurs au cent se cumulent exactement~; préfixes et catégories s'agrègent~;
les écritures erronées lèvent une exception.
\bfield{Pour aller plus loin} Enregistrer la journée simulée du tutoriel transaction par transaction et
reproduire le graphique en barres du tutoriel à partir du seul grand livre.
\end{build}

\begin{omsources}
\item Virtu Financial, \emph{Fourth Quarter 2025 Results}, communiqué de presse,
29 janvier 2026~; \emph{Fourth Quarter 2024 Results}, 29 janvier 2025.
\item Flow Traders, \emph{4Q and FY 2025 Results}, 12 février 2026~;
\emph{4Q and FY 2024 Results}, 13 février 2025.
\item Boston Consulting Group, \emph{Global Asset Management Report 2026};
\emph{Global Asset Management Report 2025}.
\item «~Alfred Winslow Jones~», \emph{Institutional Investor}, Hall of Fame
profil.
\item L.~Harris, \emph{Trading and Exchanges}, Oxford University Press, 2003,
chapitres 3, 13 et 14 --- la taxonomie classique des traders et des
revenus des dealers.
\end{omsources}

\section{Exercices}

\begin{exercise}[$\star$]\label{exo:m1:what-a-trading-firm-does:1}
Une action est cotée $24.98$ au cours acheteur et $25.02$ au cours vendeur. Donner la
\omterm{def:m1:what-a-trading-firm-does:spread}{fourchette}, le \omterm{def:m1:what-a-trading-firm-does:spread}{prix milieu}, la \omterm{def:m1:what-a-trading-firm-does:spread}{fourchette} exprimée en points de base du \omterm{def:m1:what-a-trading-firm-does:spread}{prix milieu}, et le
coût, par rapport au \omterm{def:m1:what-a-trading-firm-does:spread}{prix milieu}, d'un ordre au marché portant sur 300 actions.
\end{exercise}

\begin{exercise}[$\star$]\label{exo:m1:what-a-trading-firm-does:2}
Un \omterm{def:m1:what-a-trading-firm-does:market-maker}{teneur de marché} négocie 2 millions d'actions par jour avec une \omterm{prop:m1:what-a-trading-firm-does:capture}{captation nette}
de 0.1 cent par action. Combien gagne-t-il en une journée, et en une année de 252 jours de négociation~?
\end{exercise}

\begin{exercise}[$\star$]\label{exo:m1:what-a-trading-firm-does:3}
Pour chaque activité, dire si la société agit en \omterm{def:m1:what-a-trading-firm-does:principal-agent}{principal} ou en \omterm{def:m1:what-a-trading-firm-does:principal-agent}{agent}, et
laquelle des cinq sources de revenus de la
\cref{met:m1:what-a-trading-firm-does:sources} rémunère la société~: (a) un courtier achemine l'ordre
d'un client vers une bourse contre une commission~; (b) un dealer achète un bloc
à un client avec une décote~; (c) un gestionnaire gère le portefeuille actions d'un fonds de
retraite pour 0.15~\% par an~; (d) une société cote en permanence les deux côtés
d'un fonds indiciel coté (ETF) toute la journée~; (e) un fonds achète une action parce que son
modèle prévoit une hausse.
\end{exercise}

\begin{exercise}[$\star\star$]\label{exo:m1:what-a-trading-firm-does:4}
Un \omterm{def:m1:what-a-trading-firm-does:hedge-fund}{hedge fund} débute l'année avec 500~millions de dollars et gagne 12~\% avant frais.
Il prélève des frais de gestion de 2~\% sur les actifs d'ouverture et 20~\% du bénéfice
restant après ces frais. Calculer les deux montants de frais, le rendement net des investisseurs et
la part du bénéfice brut revenant au gérant.
\end{exercise}

\begin{exercise}[$\star\star$]\label{exo:m1:what-a-trading-firm-does:5}
Un \omterm{def:m1:what-a-trading-firm-does:market-maker}{teneur de marché} gagne une \omterm{def:m1:what-a-trading-firm-does:spread}{demi-fourchette} de 1 cent~; 14~\% de ses exécutions sont
informées et suivies d'un mouvement défavorable de 5 cents~; il reçoit un rabais de
0.20 cent et paie 0.05 cent de compensation par action. Calculer sa \omterm{prop:m1:what-a-trading-firm-does:capture}{captation nette}.
Ses coûts fixes sont de 90\,000~dollars par jour~: déterminer son volume de seuil de rentabilité.
\end{exercise}

\begin{exercise}[$\star\star$]\label{exo:m1:what-a-trading-firm-does:6}
Utiliser les chiffres 2024 de
l'\cref{dat:m1:what-a-trading-firm-does:listed}. (a) Exprimer le revenu net de trading de Flow Traders
en points de base de la valeur des produits cotés qu'elle a
négociés. Pourquoi s'agit-il d'une borne supérieure de sa captation sur ces produits~?
(b) Calculer le revenu net de trading par salarié. (c) Pour un produit à 50~euros, combien de
cents représente la captation de (a)~?
\end{exercise}

\begin{exercise}[$\star\star\star$]\label{exo:m1:what-a-trading-firm-does:7}
\emph{Programmation.} Avec le simulateur du tutoriel, exécuter \texttt{simulate\_day(7)}.
(a) Indiquer la \omterm{def:m1:what-a-trading-firm-does:spread}{fourchette} gagnée par le \omterm{def:m1:what-a-trading-firm-does:market-maker}{teneur de marché}, le P\&L de position et le
résultat net. (b) Le P\&L de position s'éloigne fortement de son espérance de $-\$9\,000$ sur de nombreuses journées.
Sans rien exécuter, expliquer quel terme du modèle crée cette
dispersion, et comment elle croîtrait si chaque ordre portait sur 400 actions et
s'il y en avait 5\,000.
\end{exercise}

\begin{exercise}[$\star\star\star$]\label{exo:m1:what-a-trading-firm-does:8}
\emph{Trouver la faille.} Un dossier de présentation affirme~: «~Notre stratégie rapporte la
\omterm{def:m1:what-a-trading-firm-does:spread}{fourchette}. La \omterm{def:m1:what-a-trading-firm-does:spread}{fourchette} moyenne de notre univers est de 6~points de base, nous négocions notre
capital dix fois par jour, donc nous gagnons 60~points de base par jour avant coûts. Nous
exécutons au moyen d'ordres au marché pour garantir nos exécutions.~» Trouver deux erreurs
indépendantes, et dire quel sera vraisemblablement le chiffre quotidien.
\end{exercise}

\section{Problème~: une semaine chez Quad Street Markets}

\begin{problem}[{Problème du week-end --- une semaine chez un petit teneur de marché}]\label{pb:m1:what-a-trading-firm-does:1}
Quad Street Markets est une \omterm{def:m1:what-a-trading-firm-does:prop}{société de trading pour compte propre} de quatorze personnes qui fait office de \omterm{def:m1:what-a-trading-firm-does:market-maker}{teneur de marché} sur
quarante grandes actions. Son propriétaire pose une seule question~: «~combien devons-nous
négocier pour survivre~?~»

\medskip\noindent\textbf{Partie I --- Ce que rapporte une action.} Une action moyenne
se négocie à 40~\$ avec une \omterm{def:m1:what-a-trading-firm-does:spread}{fourchette} de 2~cents. Quad Street négocie 6~millions d'actions
par jour, toujours passivement.
\begin{enumerate}
\item Donner la \omterm{def:m1:what-a-trading-firm-does:spread}{demi-fourchette} en cents et en points de base du cours.
\item 12~\% des exécutions de Quad Street sont informées et suivies d'un
mouvement défavorable de 5~cents. Quel est le coût de la \omterm{def:m1:what-a-trading-firm-does:adverse-selection}{sélection adverse} par action~?
\item Les bourses versent une remise de 0.20~cent~; la compensation coûte 0.02~cent.
Calculer la \omterm{prop:m1:what-a-trading-firm-does:capture}{captation nette}.
\item Calculer le revenu de négociation quotidien espéré.
\item Et le revenu annuel espéré, sur 252 jours.
\end{enumerate}

\medskip\noindent\textbf{Partie II --- Ce que coûte l'entreprise.}
\begin{enumerate}[resume]
\item Quatorze personnes coûtent 350\,000~\$ chacune par an, tout compris~; la technologie,
la colocation et les données coûtent 2.1~millions de dollars~; tout le reste coûte 0.56~million de dollars. Donner
le coût annuel et le coût par jour de négociation.
\item Donner le bénéfice quotidien et annuel espéré.
\item Donner le ratio coûts/revenus.
\end{enumerate}

\medskip\noindent\textbf{Partie III --- Atteindre le seuil de rentabilité.}
\begin{enumerate}[resume]
\item Trouver le volume quotidien de seuil de rentabilité $V^\star$.
\item Un nouveau fonds agressif devient actif sur les actions de Quad Street et la
part des exécutions informées de ce fonds monte à 16~\%. Trouver la nouvelle \omterm{prop:m1:what-a-trading-firm-does:capture}{captation nette} et le
nouveau $V^\star$.
\item À 6~millions d'actions par jour, quel est désormais le résultat quotidien~?
\item Quelle \omterm{def:m1:what-a-trading-firm-does:spread}{demi-fourchette} rétablirait la \omterm{prop:m1:what-a-trading-firm-does:capture}{captation nette} initiale~? Le pas de cotation minimal
est de 1~cent~: que peut réellement faire Quad Street~?
\item Indépendamment de la question 10, les bourses réduisent la remise à 0.15
cent. Trouver $V^\star$.
\end{enumerate}

\medskip\noindent\textbf{Partie IV --- Risque.} Revenir aux chiffres initiaux.
Le résultat quotidien a un écart-type de 22\,000~\$ et les jours sont
indépendants et approximativement normaux.
\begin{enumerate}[resume]
\item Quelle est la probabilité qu'une journée soit perdante~?
\item Calculer le ratio de Sharpe annualisé (moyenne divisée par l'écart-type du
résultat quotidien, multiplié par $\sqrt{252}$).
\item Quelle est la probabilité qu'une année soit perdante~?
\item Les propriétaires gardent 5~millions de dollars de capital dans la société. Quel est le
rendement annuel espéré de ce capital~?
\item Les volumes doublent pendant un trimestre volatil et les coûts ne bougent pas. De quel
facteur le profit quotidien augmente-t-il~?
\item Expliquer en deux phrases pourquoi les propriétaires d'un \omterm{def:m1:what-a-trading-firm-does:market-maker}{teneur de marché} apprécient les marchés volatils
alors que son responsable des risques ne les apprécie pas.
\item Donner en une phrase la réponse à la question du propriétaire, en nommant les
trois chiffres qu'il doit suivre chaque jour.
\end{enumerate}
\end{problem}

\section{Questions d'entretien}

\begin{interviewq}[$\star$ \iqroles{trader, chercheur, développeur}]\label{iq:m1:what-a-trading-firm-does:1}
Comment un \omterm{def:m1:what-a-trading-firm-does:market-maker}{teneur de marché} gagne-t-il de l'argent, et comment en perd-il~?
\end{interviewq}

\begin{interviewq}[$\star$ \iqroles{trader, chercheur, développeur, ingénieur ML}]\label{iq:m1:what-a-trading-firm-does:2}
Quelle est la différence entre une \omterm{def:m1:what-a-trading-firm-does:prop}{société de trading pour compte propre} et un \omterm{def:m1:what-a-trading-firm-does:hedge-fund}{hedge fund}~?
Pourquoi la même stratégie pourrait-elle être mise en œuvre différemment dans chacun des deux cas~?
\end{interviewq}

\begin{interviewq}[$\star\star$ \iqroles{trader, chercheur}]\label{iq:m1:what-a-trading-firm-does:3}
Vous cotez $49.99$ / $50.01$ et quelqu'un vous vend 100 actions à
$49.99$. Une minute plus tard, le marché s'établit à $49.96$ / $49.98$. Combien avez-vous
gagné ou perdu, et comment décomposeriez-vous ce chiffre en deux parties~?
\end{interviewq}

\begin{interviewq}[$\star\star$ \iqroles{chercheur, développeur}]\label{iq:m1:what-a-trading-firm-does:4}
Une société dégage un point de base net sur les 2~milliards de dollars qu'elle négocie chaque jour. Estimez
de tête son revenu annuel de négociation. Un point de base est-il une
valeur plausible~?
\end{interviewq}

\begin{interviewq}[$\star\star$ \iqroles{banque, trader}]\label{iq:m1:what-a-trading-firm-does:5}
Un client appelle pour vendre un bloc valant dix jours de volume moyen. Décrivez
deux façons dont votre desk peut aider, comment le desk est rémunéré dans chacune, et qui supporte le
risque de prix.
\end{interviewq}

\begin{interviewq}[$\star\star\star$ \iqroles{trader, chercheur}]\label{iq:m1:what-a-trading-firm-does:6}
Le mois dernier, les \omterm{def:m1:what-a-trading-firm-does:spread}{fourchettes} que vous avez cotées étaient deux fois plus larges que le mois précédent,
votre volume était le même, et vous avez gagné moins d'argent. Donnez deux explications et
la mesure qui permettrait de les distinguer.
\end{interviewq}
